\documentclass[11pt]{article}
\usepackage[margin=1in]{geometry}
\usepackage[T1]{fontenc}
\usepackage{lmodern}
\usepackage{microtype}
\usepackage{amsmath,amssymb,amsthm}
\usepackage{booktabs}
\usepackage{xcolor}
\usepackage[colorlinks=true,linkcolor=blue!55!black,urlcolor=blue!55!black]{hyperref}

\newtheorem{theorem}{Theorem}[section]
\newtheorem{lemma}[theorem]{Lemma}
\newtheorem{proposition}[theorem]{Proposition}

\newcommand{\score}{\mathcal{W}}
\newcommand{\R}{\mathbb{R}}
\newcommand{\ip}[2]{\langle #1,#2\rangle}
\newcommand{\norm}[1]{\lVert #1\rVert}

\title{\bfseries A Gram-certificate proof that\\[0.2em]
$\boldsymbol{\sigma_1(\R^2)\le 0.6934}$\\[0.6em]
\large What is formalized in \texttt{Besicovitchs-1-2}, and why the constant is not sharp}
\author{Yongxi Lin}
\date{September 2, 2026}

\begin{document}
\maketitle

\begin{abstract}
This note describes the theorem verified in the Lean~4 development
\texttt{Besicovitchs-1-2}: the planar one-dimensional rectifiability threshold satisfies
$\sigma_1(\R^2)\le 6934/10000$, improving the published bound of $0.7$. The proof combines the
six-point failure tree of the two-colour finite problem with thirty local Gram certificates. We
explain the argument, state precisely what the machine checks, and — since this is the natural
question — explain why the formalized constant is $0.6934$ rather than the sharp six-point
endpoint $s_*=0.69330642\ldots$, and what it would take to close the remaining gap of
$9.4\cdot10^{-5}$.
\end{abstract}

\section{The result}

For a Borel set $E\subset\R^2$ with $\mathcal H^1(E)<\infty$, the lower density at $x$ is
\[
\Theta^1_*(E,x)=\liminf_{r\downarrow0}\frac{\mathcal H^1(E\cap B(x,r))}{2r},
\]
and $\sigma_1(\R^2)$ is the least threshold $\beta$ such that $\Theta^1_*\ge\beta$ almost
everywhere forces countable $1$-rectifiability. Besicovitch's $1/2$ conjecture asserts
$\sigma_1(\R^2)=1/2$. The verified statement is

\begin{theorem}[\texttt{Besicovitch.sigma\_one\_plane\_le\_6934\_div\_10000}]
\[
\sigma_1(\R^2)\;\le\;\frac{6934}{10000}=0.6934 .
\]
\end{theorem}

The proof term depends on exactly the three axioms \texttt{propext}, \texttt{Classical.choice} and
\texttt{Quot.sound}. It uses no \texttt{native\_decide}, no \texttt{Lean.ofReduceBool}, no
floating-point value, no hash, and no custom axiom.

\section{Why the constant is not sharp}
\label{sec:why}

The six-point finite problem this development solves has an exact optimal constant,
\[
\theta_6=s_*=0.693306421825904872690678414403710951\ldots,
\]
an algebraic number pinned by an isolated pair of radical equations. So a natural target is
$\sigma_1(\R^2)\le s_*$, and an earlier form of this repository pursued exactly that. The reason
the formalized statement is the weaker rational bound is not that the mathematics fails at $s_*$;
it is that \emph{tightness destroys the certificates}.

At $s_*$ the extremal configuration is attained: a half-turn pair of chords at which three
structurally different packing mechanisms tie \emph{simultaneously}. The weighted score that the
argument must bound above by zero therefore has value exactly $0$ at an interior point of the
domain, with vanishing gradient. Any certificate that establishes a non-strict bound at a
degenerate maximum must resolve the geometry near that point to the precision of the degeneracy.
Concretely, in the earlier Bernstein formulation the margin available away from the extremiser was
about $10^{-8}$, and the adaptive subdivision needed to certify the neighbourhood of the tie grew
until one of the ten coordinate charts admitted no finite tree at all: on that chart the score is
identically $0$ at the extremiser, so no strict Positivstellensatz certificate of the assumed shape
exists.

Moving the chord to the rational value
\[
c=\frac{3467}{2500}=1.3868,\qquad s=\frac c2=\frac{6934}{10000},
\]
breaks every tie at once. The three mechanisms no longer agree, the extremal configuration is no
longer feasible, and the weighted score acquires a uniform negative margin of about
$5.8\cdot10^{-4}$ — roughly four orders of magnitude more room than at $s_*$. That margin is what
makes a \emph{small, fixed} certificate family sufficient: thirty rectangles suffice where the
sharp problem needed an unbounded adaptive tree.

The cost of this trade is quantified exactly: the gap between what is proved and the sharp
six-point constant is
\[
\frac{6934}{10000}-s_*\;=\;9.357817\cdot10^{-5}.
\]
The gap to the conjecture itself is of course far larger; $\sigma_1(\R^2)=1/2$ remains open, and
the six-point method cannot reach below $\theta_6$ no matter how it is certified. A companion
theorem in \texttt{Solution.lean}, \texttt{sStar\_le\_6934\_div\_10000}, records that the exact
endpoint does lie below the rational target, so the two statements are consistent and the
formalized bound is genuinely weaker only by the amount displayed above.

\medskip
There is a second, upper constraint on the choice, which is why the constant is $0.6934$ and not,
say, $0.699$. The routing separators that prune the case tree are themselves inequalities in the
chord, and they cease to hold once $c$ exceeds roughly $1.386850$. The rational chord above sits
just under that ceiling while clearing $s_*$ from above, so the interval of admissible targets is
narrow, and $6934/10000$ is a natural point inside it.

\section{Input from the six-point failure tree}

Let $e,p_1,p_2,w_1,w_2$ lie in a real inner-product space with
\[
\norm e=1,\qquad \norm{p_i}\le1,\qquad \norm{w_i}\le1,
\]
and suppose the two sibling pairs are separated,
\begin{equation}
c\le\norm{p_1-p_2},\qquad c\le\norm{w_1-w_2}.
\label{eq:sep}
\end{equation}
The finite failure tree — retained unchanged from the sharp development — reduces the failure of
all nine packings at a matched sibling endpoint to three scalar inequalities
\begin{equation}
q_1\ge0,\qquad q_2>0,\qquad q_3>0,
\label{eq:slacks}
\end{equation}
in Lean \texttt{firstActiveFailureSlack\_nonneg}, \texttt{secondActiveFailureSlack\_pos} and
\texttt{thirdActiveFailureSlack\_pos}.

The weighted pair score is, for weights $\lambda,\mu>0$,
\begin{align*}
\score_{c,\lambda,\mu}
={}&(1+\lambda)\norm{e-p_1-w_1}+\norm{e-p_2-w_2}\\
&+\tfrac\mu2\bigl(\norm{e-p_1}+\norm{e-w_1}+\norm{e-p_1-w_2}+\norm{e-w_1-p_2}\bigr)\\
&-\tfrac{(c-1)(\lambda/2+\mu)}2\bigl(\norm{p_1}+\norm{w_1}\bigr)
-\tfrac{(c+1)\lambda/2+3c\mu}{2}\bigl(\norm{p_2}+\norm{w_2}\bigr)\\
&-\Bigl(2c(2c-1)+\tfrac\lambda2(3c^2-3c+2)+\mu(c^2-c)\Bigr),
\end{align*}
and a purely algebraic identity gives
$q_1+\lambda q_2+\mu q_3=\score_{c,\lambda,\mu}$ for \emph{every} $\lambda,\mu$
(\texttt{weightedPairScore\_configuration\_eq\_activeFailureCombination}). Hence
\eqref{eq:slacks} forces $\score_{c,\lambda,\mu}>0$.

The contradiction we must produce is an upper bound $\score<0$. In the sharp development the
weights were the Cramer weights of the endpoint system, algebraic numbers of high degree. Because
the identity holds for arbitrary weights, we are free to pick
\begin{equation}
\lambda=\frac1{12},\qquad \mu=\frac{13}{14},
\label{eq:weights}
\end{equation}
which is what makes every subsequent coefficient a small rational. This substitution alone removes
the entire exact-endpoint apparatus from the dependency path of the theorem.

\section{One local Gram certificate}

Fix a rectangle of second-child radii
\begin{equation}
l_P\le\norm{p_2}\le u_P,\qquad l_W\le\norm{w_2}\le u_W .
\label{eq:box}
\end{equation}
The triangle inequality and \eqref{eq:sep} give $c-u_P\le\norm{p_1}\le1$ and
$c-u_W\le\norm{w_1}\le1$; note that the lower bound uses $u_P$, not $l_P$.

Write $v_0=e$, $v_1=p_1$, $v_2=p_2$, $v_3=w_1$, $v_4=w_2$. With the weights \eqref{eq:weights}
the six positively weighted terms of $\score$ are $\norm{\xi_t}$ for $\xi_t=\sum_j A_{tj}v_j$,
\[
A=\begin{pmatrix}
1&-1&0&-1&0\\ 1&0&-1&0&-1\\ 1&-1&0&0&0\\
1&0&0&-1&0\\ 1&-1&0&0&-1\\ 1&0&-1&-1&0
\end{pmatrix},
\qquad
\omega=\Bigl(\tfrac{13}{12},1,\tfrac{13}{28},\tfrac{13}{28},\tfrac{13}{28},\tfrac{13}{28}\Bigr).
\]
Three elementary inequalities linearize the problem. For $\alpha_t>0$ the quadratic tangent
$\omega_t\norm{\xi_t}\le\alpha_t\norm{\xi_t}^2+\omega_t^2/(4\alpha_t)$ replaces each norm by a
quadratic form; the radial secant $-dr\le-\frac d{l+u}r^2-\frac{dlu}{l+u}$ (valid for
$0\le l\le r\le u$) replaces each of the four radial penalties; and for
$\eta_P,\eta_W\ge0$ the separation multipliers $\eta_P(\norm{p_1-p_2}^2-c^2)\ge0$ and
$\eta_W(\norm{w_1-w_2}^2-c^2)\ge0$ may be added freely. Expanding the squared norms,
\begin{equation}
\score\;\le\;C+\sum_{i=0}^4 D_i\norm{v_i}^2+2\!\!\sum_{0\le i<j\le4}\!\!Q_{ij}\ip{v_i}{v_j},
\label{eq:quad}
\end{equation}
with $C,D,Q$ explicit rational functions of the certificate parameters and the box.

It remains to remove the inner products. A certificate carries a rational $3\times5$ matrix $F$;
put $B=F^{\mathsf T}F$, let $\rho_{ij}=-Q_{ij}-B_{ij}$, and let
$\varepsilon_{ij}=\operatorname{sign}\rho_{ij}$. The matrix
\begin{equation}
M=B+\sum_{i<j}|\rho_{ij}|\,(e_i+\varepsilon_{ij}e_j)(e_i+\varepsilon_{ij}e_j)^{\mathsf T}
\label{eq:psd}
\end{equation}
is a sum of rank-one positive semidefinite matrices, hence positive semidefinite, and by
construction $M_{ij}=-Q_{ij}$ off the diagonal. Since the Gram matrix
$G=(\ip{v_i}{v_j})$ is also positive semidefinite, the Schur product theorem gives
$\sum_{i,j}M_{ij}\ip{v_i}{v_j}\ge0$, that is
\[
0\;\le\;\sum_i m_i\norm{v_i}^2-2\sum_{i<j}Q_{ij}\ip{v_i}{v_j},
\qquad m_i=B_{ii}+\sum_{j\ne i}|\rho_{ij}| .
\]
Adding this to \eqref{eq:quad} cancels every inner product and leaves
$\score\le C+\sum_i\beta_i\norm{v_i}^2$ with $\beta_i=D_i+m_i$. Finally each
$\beta_i\norm{v_i}^2$ is bounded on its radius interval by
$T(\beta;l,u)=(\beta)_+u^2-(\beta)_-l^2$, giving a single rational number
\begin{equation}
U=C+\beta_0+T(\beta_1;c-u_P,1)+T(\beta_2;l_P,u_P)+T(\beta_3;c-u_W,1)+T(\beta_4;l_W,u_W).
\label{eq:U}
\end{equation}

\begin{proposition}[\texttt{weightedPairScore\_le\_of\_gramCertificate}]
If $\alpha_t>0$, $\eta_P,\eta_W\ge0$ and $U\le-\frac1{2000}$, then every configuration satisfying
\eqref{eq:sep} and \eqref{eq:box} in the unit ball has $\score\le-\frac1{2000}<0$.
\end{proposition}

No eigenvalue computation, numerical positive-definiteness test, or polynomial arithmetic occurs
anywhere: positive semidefiniteness is established \emph{identically} from the rank-one
decomposition \eqref{eq:psd}.

\section{The finite cover}

Every feasible second-child radius lies in $[c-1,1]$. Seven intervals cover that range,
\[
\begin{array}{lll}
I_0=[c-1,\tfrac12], & I_1=[\tfrac12,\tfrac35], & I_2=[\tfrac35,\tfrac7{10}],\\[2pt]
I_3=[\tfrac7{10},\tfrac34], & I_4=[\tfrac34,\tfrac45], & I_5=[\tfrac45,\tfrac9{10}],\\[2pt]
I_6=[\tfrac9{10},1]. &&
\end{array}
\]
The score is symmetric under exchanging the two sibling pairs, so only the $28$ ordered pairs
$I_i\times I_j$ with $i\le j$ are needed. The single square $I_2\times I_2$ carries too little
margin and is split once more using $I_{2a}=[\tfrac35,\tfrac{13}{20}]$ and
$I_{2b}=[\tfrac{13}{20},\tfrac7{10}]$, contributing three rectangles instead of one. The total is
$28-1+3=30$.

\begin{theorem}[\texttt{weightedGeometricBound\_gram}]
For every $e,p_1,p_2,w_1,w_2$ in a real inner-product space with $\norm e=1$,
$\norm{p_i},\norm{w_i}\le1$ and \eqref{eq:sep},
\[
\score_{c,1/12,13/14}\;\le\;-\frac1{2000}.
\]
\end{theorem}

Combined with the failure tree this contradicts $\score>0$, so the six-point finite property holds
at $s=6934/10000$; the already formalized six-point-to-BPC and BPC-to-rectifiability transfers
then give the theorem.

The binding certificate is the one on $I_3\times I_5$, whose exact upper bound is
\[
U=-\frac{141718938311359411043}{245112420775950000000000}=-0.000578179\ldots<-\frac1{2000},
\]
so the whole family clears the uniform threshold with about $16\%$ to spare.

\section{The lower bound}
\label{sec:lower}

The registry review observed that Lean's real infimum returns $0$ on the empty set, so an upper
bound on $\sigma_1$ could in principle hold vacuously. The formalization therefore also proves
the classical lower bound
\[
\tfrac12\le\sigma_1(\R^2),
\]
which pins the threshold to $[\tfrac12,0.6934]$ and excludes the vacuous case outright. The
statement \texttt{one\_half\_le\_sigma\_one\_plane} is compared alongside the upper bound.

A lower bound needs a set of finite positive length that is not countably rectifiable and yet
has lower density at least $\tfrac12$ almost everywhere. Besicovitch proposed one in 1938 and
Dickinson proved its properties in 1939; we follow Capdevila's 2026 account. Let
\[
\alpha_n=2^{-n^2},\qquad \beta_n=\alpha_n/n,\qquad
g=\sum_{n\ge1}f_n,
\]
where $f_n$ is the square wave of period $2\alpha_n$ and amplitude $\beta_n$, and let $\Pi$ be
the graph of $g$ over $[0,1]$. Two estimates carry everything. Inside a level-$n$ cell
$[i\alpha_n,(i+1)\alpha_n)$ the function varies by at most $4\alpha_{n+1}/(n+1)$; across a
level-$n$ cell boundary it jumps by at least $\alpha_n/n$. Mathlib contains no notion of
rectifiability and no purely unrectifiable set, so the example and its three properties are
built from scratch.

\emph{Measure.} The first-coordinate projection is $1$-Lipschitz, so $\mathcal H^1(\Pi(A))\ge|A|$
for every $A\subset[0,1]$; covering the graph over a cell by a cylinder of diameter at most
$2\alpha_n$ gives $\mathcal H^1(\Pi(A))\le 2|A|$, first on intervals, then on open sets through
the cells they contain, then on arbitrary sets by outer regularity. The sharp constant $1$ is
never needed. In particular the graph over a null set is null.

\emph{Density.} At a point $x\in(0,1)$ and a small radius $r$, choose the level $n$ with
$\alpha_{n+1}<r\le\alpha_n$. The graph over the part of $x$'s cell within $\theta r$ of $x$ lies
in the ball of radius $r$ about $(x,g(x))$, because the vertical variation inside the cell is
$O(r/n)$; and that part has length at least $\theta r$, since one side of $x$ reaches at least
that far within a cell of length $\ge r$. Hence the density is at least $\theta/2$ for every
$\theta<1$, so at least $\tfrac12$.

\emph{Unrectifiability.} This is the step that is elementary here and not in general. If a
Lipschitz curve met $\Pi$ in positive measure, Rademacher's theorem and Mathlib's partition of a
differentiable map into pieces on which it approximates an invertible linear map would produce
a set $A\subset[0,1]$ of positive measure on which $g$ is Lipschitz. But $g$ jumps by at least
$\alpha_n/n$ across every level-$n$ boundary, so a set on which $g$ is $L$-Lipschitz cannot have
points within $\alpha_n/(2n(L+1))$ of a boundary on both sides. That alone does not make $A$
null — a single such hole is a vanishing fraction of any ball — but a Lebesgue density point of
$A$ cannot lie that close to a level-$n$ boundary for infinitely many $n$, because the ball of
twice that radius about it would then carry a hole of relative size $\tfrac14$. So almost every
point of $A$ eventually keeps its distance from the grid on \emph{both} sides, and the set of
points that do so at all fine levels is null by a nested recursion: inside each cell the
surviving set is an interval, the grid points of the next level carve disjoint gaps out of it
whose total is a fraction $1/(n(L+1))$ of its length, and $\sum 1/n=\infty$. No projection
theorem, no tangent theory and no area formula are used.

\section{What the machine checks}

The certificate parameters were found by an unverified numerical search and are stored as exact
rationals with denominator $10^4$. Nothing about how they were found enters the proof: Lean
recomputes $C$, $D$, $Q$, $B$, $\rho$, $m$, $\beta$ and $U$ from the stored data by exact rational
arithmetic and checks $\alpha_t>0$, $\eta_P,\eta_W\ge0$ and $U\le-1/2000$ for each of the thirty
records (\texttt{gramCertificates\_valid}). The positive semidefiniteness is a theorem, not a
computation. The cover is checked by deciding, for each of the $64$ ordered pairs of radius bands,
that the band rectangle sits inside the rectangle of the assigned certificate.

\begin{center}
\begin{tabular}{@{}lr@{}}
\toprule
stage & cost\\
\midrule
thirty certificate validity checks & $43$ s\\
band cover & $3$ s\\
Gram certificate core & $7$ s\\
\bottomrule
\end{tabular}
\end{center}

For comparison, the six-variable Bernstein formulation this replaces required an estimated
$29$ CPU-hours of kernel reduction over $3072$ adaptive leaves, and was in any case incomplete:
one of the ten coordinate charts had no certificate tree.

\section{Provenance of the certificate data}

The thirty certificates were checked three times independently before formalization: by an exact
rational verifier, by a symbolic reconstruction using a computer algebra system, and by a third
transcription written directly against the definitions above, which reproduced all thirty
published values of $U$ bit-for-bit. That agreement also validates the transcription into Lean,
since a single mistyped integer would change $U$.

\end{document}
